%% 第十章 总结与展望
\chapter{总结与展望}

\section{全文总结}

本文面向低空反无人机场景，以单目可见光相机为唯一传感器，设计并实现了一套从检测、跟踪到定位与显示的完整系统（DRLS）。围绕小目标检测难、遮挡跟踪易中断、仿真与真实域差异大三个核心问题，本文完成了以下工作。

（1）\textbf{小目标检测模型改进}。以 YOLOv8n 为基础，增加 stride=4 的 P2 高分辨率检测头与 EMA 高效多尺度注意力，构建 YOLOv8n-P2-EMA 模型。在 Det-Fly 数据集上按序列划分训练/验证集，避免数据泄漏。四组消融实验表明，本文模型以 2.94\,M 参数量取得 mAP50 83.5\%、Recall 79.2\%、AP\_small 9.8\%，相对基线分别提升 2.2、3.3 与 3.7 个百分点，AP\_small 相对提升 60.7\%，推理速度 98.4\,fps。本文进一步解释了 P2 与 EMA 的互补机制：P2 提供"看清"小目标的分辨率，EMA 提供"认准"小目标的判别力，二者联合才能发挥最佳效果。

（2）\textbf{数据工程}。完成 Det-Fly 的格式转换、统计分析与序列级划分；自建 AprilTag 数据集，使用 9$\times$6、25\,mm 棋盘格完成相机标定，为定位精度评估提供可信真值；构建 450 帧仿真自动标注数据集，利用护航相机注视特性实现零人工标注，并按距离分层覆盖死区两端。

（3）\textbf{三态跟踪状态机}。设计 SEARCH/LOCKED/COAST 三种状态，采用 0.25 维持门限与 0.55 初锁门限的双门限策略抑制误锁，采用 140\,px 空间门限的最近邻关联抑制身份跳变，采用 1.0\,s 卡尔曼预测维持抑制短时遮挡。参数选择均有物理依据，算法复杂度 $O(N)$。实验表明，TCP 遮挡场景的航迹保持率达到 88.9\%--94.5\%，真实场景录像达到 99.3\%。

（4）\textbf{单目定位与自适应滤波}。由检测框中心直接解算方位角与俯仰角，由框高与等效尺寸解算距离，推导了测距误差随距离平方增长的规律；设计六维状态的自适应卡尔曼滤波器，采用各向异性量测噪声与 Sage-Husa 在线估计。仿真闭环中，单帧测距 RMSE 0.26\,m、AKF 稳态 MedAE 0.13\,m、1\,s 预测 ADE 中位 0.26\,m、方位角误差中位 0.05\textdegree；实拍视频中，AKF 相对单帧测距 RMSE 降低 32\%、抖动降低 33\%。

（5）\textbf{仿真闭环与系统集成}。改造 PX4+jMAVSim 闭环，新增脚本目标机、护航相机、真值广播、离屏帧流与城市场景；通过 2228 帧的距离分层统计与真值几何审计，定量证伪"俯视导致失效"的假设，将检测死区归因于尺度分布偏移；提出仿真域微调方法，用 450 帧分层数据将 3--10\,m 检出率从约 10\% 提升至 100\%。基于 PySide6 实现三页签监视大屏，采用双进程架构与 UDP/JSON 通信，完成态势地图、光电画面、仪表告警、配置持久化与 14 用例自动化测试。

\section{主要创新点}

本文的主要创新点可以归纳为以下四点。

（1）\textbf{面向无人机小目标的 P2+EMA 互补结构及其机理解释}。不仅通过消融实验验证了改进的有效性，还解释了为什么 EMA 单独使用会失效、与 P2 联合后却有效，为小目标检测的网络设计提供了可复用的分析框架。

（2）\textbf{基于真值几何审计的检测死区机理分析方法}。面对"模型在仿真中失效"的现象，本文没有停留在"渲染不真实"的定性解释，而是通过俯仰角分布对比证伪视角假设，通过距离分层统计锁定尺度因素，并用微调实验闭环验证。这种"现象—假设—证伪—验证"的路径具有方法论价值。

（3）\textbf{零人工标注的仿真域微调数据构建方法}。利用护航相机始终注视目标的特性，将框中心约束在光轴附近，由真值距离反算框高，实现了 450 帧数据的自动标注。该方法成本极低、可扩展性强，适用于任何具备真值广播与可控相机的仿真环境。

（4）\textbf{面向反无人机监视的完整工程实现}。将检测、三态跟踪、单目定位、地图态势、光电显示、告警与测试整合为可运行的双进程系统，并在统一口径下完成端到端评估。系统不依赖专用硬件，便于在园区、楼宇与临时活动现场部署。

\section{不足与展望}

\subsection{存在的不足}

（1）\textbf{远距离能力有限}。实拍视频在 10--14\,m 的检出率仅 5.6\%，说明当前模型与 720p 分辨率组合无法覆盖更远的告警距离。这是单目视觉的物理限制，也是系统走向实用的主要瓶颈。

（2）\textbf{相机标定不完整}。实拍评估中相机未做严格标定，焦距按视场角估算，等效尺寸借用仿真值，导致绝对测距误差达到米级。虽然相对比较有效，但绝对精度不足以支撑实际指挥。

（3）\textbf{单目标假设}。系统面向单目标场景，未处理多架无人机同时出现的情况。当画面中出现多个目标时，状态机可能在不同目标之间跳变。

（4）\textbf{仿真场景简化}。仿真场景的建筑、树木与光照仍与真实世界有差距，仿真结论不能直接外推到外场；同时，仿真中的遮挡主要由静态建筑造成，缺少动态遮挡（如飞鸟、飘动的树叶）。

（5）\textbf{测试集偏差}。退化实验的测试帧由 simft 模型筛选，对 simft 有利，三组模型的比较不完全公平。

\subsection{未来工作}

（1）\textbf{提升远距离能力}。可以从三个方向入手：使用更高分辨率输入（如 1280 或 1536）并配合切片推理；使用长焦镜头与可变速云台；引入双目视觉或激光测距辅助深度估计。

（2）\textbf{完善相机标定与尺寸自适应}。使用棋盘格完成实拍相机标定，替代视场角估算；根据目标机型选择等效尺寸，或利用多帧几何一致性在线估计尺寸。

（3）\textbf{扩展多目标跟踪}。在现有状态机之上增加轨迹池与数据关联层，借鉴 SORT/DeepSORT 的思想，实现多目标的同时跟踪与威胁排序。

（4）\textbf{增强域适应能力}。探索域随机化、风格迁移与自监督预训练，减少对目标域标注的依赖；同时构建更真实的仿真场景，缩小 sim-to-real gap。

（5）\textbf{多传感器融合}。将视觉与射频、声学、雷达结合，利用射频进行远距离预警、利用视觉进行近程确认与定位，形成互补的探测体系。

（6）\textbf{外场试验}。在合规空域开展真实无人机飞行试验，采集多场景、多机型、多光照的数据，验证系统的实际可用性，并根据外场反馈迭代参数与模型。

\section{结语}

反无人机是一个涉及探测、识别、跟踪、定位与处置的系统工程，单目视觉只是其中的一个环节。本文的工作表明：通过针对性的模型改进、严谨的机理分析与完整的工程实现，单目视觉系统可以在近程反无人机场景中发挥重要作用；但同时也要清醒地认识到，单目视觉的物理限制决定了它无法单独承担全部探测任务。未来的反无人机系统必然是多传感器、多算法、多层级协同的体系，本文的工作可以作为其中的一个可复现、可扩展的视觉子系统。
